\documentclass[12pt,xcolor=dvipsnames]{article}
\usepackage[utf8]{inputenc}
\usepackage[vietnamese]{babel}
\usepackage{amstext,latexsym,amsbsy,amssymb,amsmath,amsthm,amsfonts,exscale,eucal}
\usepackage{appendix}
\usepackage{algorithm}
\usepackage{graphicx}
\usepackage{xcolor}
\usepackage{hyperref}
\usepackage{indentfirst}
\usepackage{enumitem}
\usepackage{diagbox}

\usepackage[left=2.00cm, right=2.00cm, top=2.00cm, bottom=2.00cm]{geometry}
\setlength{\parskip}{5pt}%khoảng cách dòng
\renewcommand{\baselinestretch}{1.15}
\theoremstyle{definition}
\newtheorem{vd}{Ví dụ}[section]
\newtheorem*{giai}{Giải}
\newtheorem{bt}{Bài tập}
\renewcommand{\l}{\left}
\renewcommand{\r}{\right}
\date{}

\begin{document}
\begin{bt}
Viết hàm thực hiện hồi quy đa thức (local polynomial regression), với input $p$ là số nguyên có các giá trị $0, 1, 2, \ldots$ Trong đó, băng thông $h$ được truyền vào là 1 số hoặc không truyền vào (đặt mặc định $h = \texttt{NULL}$), nếu không truyền vào thì trong thân hàm cần có một đoạn tính $h$ theo CV hoặc GCV.
\end{bt}

\begin{giai}
\textbf{1. Ý tưởng thuật toán}

Hồi quy đa thức cục bộ (Local Polynomial Regression) ước lượng hàm hồi quy $m(x) = \mathbb{E}[Y|X=x]$ tại mỗi điểm $x_0$ bằng cách khớp một đa thức bậc $p$ với trọng số kernel:
\begin{align}
\min_{\beta_0,\beta_1,\ldots,\beta_p} \sum_{i=1}^n \left(Y_i - \beta_0 - \beta_1(X_i-x_0) - \cdots - \beta_p(X_i-x_0)^p\right)^2 K\left(\frac{X_i-x_0}{h}\right),
\end{align}
trong đó $K(\cdot)$ là hàm kernel (ở đây sử dụng kernel Gaussian), và $h>0$ là băng thông.

Nghiệm của bài toán trên cho ta ước lượng $\hat{m}(x_0) = \hat{\beta}_0$, tức là hệ số tự do của đa thức cục bộ.

\textbf{2. Cấu trúc hàm} 

Hàm \texttt{locpoly\_reg} có cấu trúc như sau:
\begin{verbatim}
locpoly_reg(x, y, p, h = NULL, eval_x = NULL)
\end{verbatim}
với các tham số:
\begin{itemize}
    \item \texttt{x}: vector các giá trị biến độc lập,
    \item \texttt{y}: vector các giá trị biến phụ thuộc,
    \item \texttt{p}: bậc đa thức ($p=0,1,2,\ldots$),
    \item \texttt{h}: băng thông (nếu \texttt{NULL} thì tự động chọn qua CV),
    \item \texttt{eval\_x}: các điểm cần ước lượng (mặc định là \texttt{x}).
\end{itemize}

\textbf{3. Chọn băng thông tự động qua Cross-Validation}

Khi $h=\texttt{NULL}$, ta sử dụng Leave-One-Out Cross-Validation (LOOCV) để chọn băng thông tối ưu:
\begin{align}
\text{CV}(h) = \frac{1}{n}\sum_{i=1}^n \left(Y_i - \hat{m}_{-i}(X_i;h)\right)^2,
\end{align}
trong đó $\hat{m}_{-i}(X_i;h)$ là ước lượng tại $X_i$ khi loại bỏ điểm thứ $i$ khỏi tập huấn luyện.

Do việc tối ưu hóa liên tục tốn kém, ta sử dụng grid search trên 8 giá trị băng thông phân bố logarithmic trong khoảng $[h_{\min}, h_{\max}]$:
\begin{align}
h_{\min} = \max\left\{2\cdot\min_{i\ne j}|X_i-X_j|, \frac{\text{range}(X)}{100}\right\}, \quad
h_{\max} = \frac{\text{range}(X)}{2}.
\end{align}

\textbf{4. Triển khai}

Code R được lưu tại:
\begin{itemize}
    \item \texttt{Exercise/Exercise 1/exercise\_03.R}: Hàm chính \texttt{locpoly\_reg}.
    \item \texttt{Exercise/Exercise 1/test\_exercise\_03.R}: Script kiểm thử và vẽ đồ thị.
\end{itemize}

Biểu đồ kết quả được lưu tại \texttt{Exercise/Exercise 1/exercise\_03\_result.png}.
\end{giai}
\end{document}